% paper/sections/09_full_algorithm.tex
\section{完全アルゴリズムフロー：離散形式}
\label{sec:algorithm}

時刻 $t^n$ から $t^{n+1} = t^n + \Delta t$ への計算手順を，離散化された形式で以下にまとめる．
添字 $(i,j)$ は格子点のインデックスを，上付き文字 $n, *, n+1$ は時間ステップを表す．
以下の離散化オペレータを用いる．

\begin{tcolorbox}[defbox, title={離散化オペレータの定義}]
\begin{itemize}
  \item $D_x^{(1)}, D_y^{(1)}$：CCD 法による1次微分 $(\partial/\partial x,\, \partial/\partial y)$
  \item $D_x^{(2)}, D_y^{(2)}$：CCD 法による2次微分 $(\partial^2/\partial x^2,\, \partial^2/\partial y^2)$
  \item $D_{xy}^{(11)}$：CCD 法による混合微分 $\partial^2/\partial x\partial y$
  \item $\mathcal{C}_{\text{WENO}}(\bu,\, q)$：速度 $\bu$ による WENO5 上流移流オペレータ
          $\displaystyle= u\,D_x^{(1)} q + v\,D_y^{(1)} q$ （保存形）
  \item $\mathcal{L}_{\text{visc}}(\mu,\, \bu)$：CCD 法による粘性拡散オペレータ
          $\displaystyle= D_x^{(1)}\!\left[\mu(D_x^{(1)} u + D_y^{(1)} v)\right] + D_y^{(1)}\!\left[\mu(D_x^{(1)} v + D_y^{(1)} v)\right]$ （各速度成分に適用）
  \item $\text{Div}_h(\bm{q})$：Rhie-Chow 補間を内包した離散発散 $D_x^{(1)} q_x + D_y^{(1)} q_y$
  \item $\text{Grad}_h(p)$：$(D_x^{(1)} p,\; D_y^{(1)} p)$
\end{itemize}
\end{tcolorbox}

\begin{tcolorbox}[algbox, title={時刻 $t^n \to t^{n+1}$ の離散計算フロー}]

\textbf{1. CLS 移流 (Advection)} \\
保存形レベルセット変数 $\psi^n$ を，速度場 $\bm{u}^n$ を用いて時間発展させる．
ここでは3次のルンゲ=クッタ法（TVD-RK3）を用いる．
\begin{align*}
  \psi^{(1)} &= \psi^n - \Delta t \, \text{Div}_h(\psi^n \bm{u}^n) \\
  \psi^{(2)} &= \frac{3}{4}\psi^n + \frac{1}{4}\psi^{(1)} - \frac{1}{4}\Delta t \, \text{Div}_h(\psi^{(1)} \bm{u}^n) \\
  \psi^* &= \frac{1}{3}\psi^n + \frac{2}{3}\psi^{(2)} - \frac{2}{3}\Delta t \, \text{Div}_h(\psi^{(2)} \bm{u}^n)
\end{align*}
ここで $\text{Div}_h(\cdot)$ は5次精度WENOスキームで計算された発散オペレータを表す．

\textbf{2. 再初期化 (Reinitialization)} \\
移流によって鈍化した界面プロファイル $\psi^*$ を整形する．
式 \eqref{eq:cls_reinit_iso} で表される偏微分方程式を，仮想時間 $\tau$ について複数ステップ解く．
\[ \psi^{n+1} = \text{Reinitialize}(\psi^*) \]

\textbf{3. グリッド・物性更新 (Update Properties)} \\
更新された $\psi^{n+1}_{i,j}$ に基づき，各格子点での距離関数 $\phi$ および物理的物性値を再計算する．
\begin{align*}
  \phi^{n+1}_{i,j} &= H_\varepsilon^{-1}(\psi^{n+1}_{i,j}) \quad (\text{Newton法で解く}) \\
  \rho^{n+1}_{i,j} &= \rho(\phi^{n+1}_{i,j}) \\
  \mu^{n+1}_{i,j} &= \mu(\phi^{n+1}_{i,j})
\end{align*}

\textbf{4. 曲率計算 (Curvature)} \\
$\phi^{n+1}$ の空間微分をCCD法で高精度に計算し，式 \eqref{eq:curvature_2d} に直接代入して曲率を求める．
\begin{align*}
 (\phi_x)_{i,j} &= D_x^{(1)} \phi^{n+1}_{i,j}, \quad
 (\phi_y)_{i,j} = D_y^{(1)} \phi^{n+1}_{i,j} \\
 (\phi_{xx})_{i,j} &= D_x^{(2)} \phi^{n+1}_{i,j}, \quad
 (\phi_{yy})_{i,j} = D_y^{(2)} \phi^{n+1}_{i,j}, \quad
 (\phi_{xy})_{i,j} = D_{xy}^{(11)} \phi^{n+1}_{i,j}
\end{align*}
\begin{equation*}
 \kappa^{n+1}_{i,j} = -\frac{(\phi_y)_{i,j}^2 (\phi_{xx})_{i,j}
                              - 2(\phi_x)_{i,j}(\phi_y)_{i,j}(\phi_{xy})_{i,j}
                              + (\phi_x)_{i,j}^2 (\phi_{yy})_{i,j}}
                           {\max\!\bigl((\phi_x)_{i,j}^2 + (\phi_y)_{i,j}^2,\;
                            \epsilon_{\text{norm}}^2\bigr)^{3/2}}
\end{equation*}

\textbf{5. 運動量予測 (Predictor Step)} \\
$x$成分と$y$成分それぞれについて，格子点 $(i,j)$ での離散式を以下に示す．

\underline{$x$成分（$u^*_{i,j}$ を求める）：}
\begin{align*}
  \frac{\rho^{n+1}_{i,j}}{\Delta t}\bigl(u^*_{i,j} - u^n_{i,j}\bigr)
  &= -\rho^{n+1}_{i,j}\Bigl[u^n_{i,j}\,D_x^{(1)} u^n_{i,j} + v^n_{i,j}\,D_y^{(1)} u^n_{i,j}\Bigr] \\
  &\quad + \frac{1}{2}\Bigl[D_x^{(1)}\!\bigl(\mu^{n+1}(D_x^{(1)} u^* + D_y^{(1)} v^*)\bigr)
                          + D_y^{(1)}\!\bigl(\mu^{n+1}(D_y^{(1)} u^* + D_x^{(1)} v^*)\bigr) \\
  &\qquad\qquad    + D_x^{(1)}\!\bigl(\mu^{n+1}(D_x^{(1)} u^n + D_y^{(1)} v^n)\bigr)
                          + D_y^{(1)}\!\bigl(\mu^{n+1}(D_y^{(1)} u^n + D_x^{(1)} v^n)\bigr)\Bigr]_{i,j} \\
  &\quad + \frac{\kappa^{n+1}_{i,j}\,(D_x^{(1)}\psi^{n+1})_{i,j}}{We}
\end{align*}

\underline{$y$成分（$v^*_{i,j}$ を求める）：}
\begin{align*}
  \frac{\rho^{n+1}_{i,j}}{\Delta t}\bigl(v^*_{i,j} - v^n_{i,j}\bigr)
  &= -\rho^{n+1}_{i,j}\Bigl[u^n_{i,j}\,D_x^{(1)} v^n_{i,j} + v^n_{i,j}\,D_y^{(1)} v^n_{i,j}\Bigr] \\
  &\quad + \frac{1}{2}\Bigl[\text{（$x$成分と対称，$u \leftrightarrow v$，$x \leftrightarrow y$ と置換）}\Bigr]_{i,j} \\
  &\quad - \frac{\rho^{n+1}_{i,j}}{Fr^2}
         + \frac{\kappa^{n+1}_{i,j}\,(D_y^{(1)}\psi^{n+1})_{i,j}}{We}
\end{align*}

移流項（第1行）は陽解法（時刻 $t^n$ の値のみ），粘性項（第2--3行）は Crank--Nicolson（$\bm{u}^*$ と $\bm{u}^n$ の平均）で離散化する．
$\bm{u}^*$ を含む粘性項のため，全格子点をまとめると $\bm{u}^* = (u^*, v^*)$ に関する\textbf{ブロック三重対角線形方程式系}が得られ，これをガウス消去法等で解く．

\textbf{6. 圧力 Poisson 方程式の求解 (Pressure Solve)} \\
次ステップの速度場 $\bm{u}^{n+1}$ が非圧縮条件 $\text{Div}_h(\bm{u}^{n+1}) = 0$ を満たすように，圧力 $p^{n+1}$ を求める．
これは以下の形式の離散的な Poisson 方程式となる．
\[ \text{Div}_h\left( \frac{1}{\rho^{n+1}} \text{Grad}_h(p^{n+1}) \right)_{i,j} = \frac{\text{Div}_h(\bm{u}^*_{i,j})}{\Delta t} \]
ここで，$\text{Div}_h, \text{Grad}_h$ は Rhie-Chow 補間を内包した離散的な発散および勾配オペレータである．
この大規模な疎行列方程式を反復法ソルバー（例：BiCGSTAB）で解き，$p^{n+1}$ を得る．

\textbf{7. 速度補正 (Corrector Step)} \\
得られた圧力勾配を用いて，予測速度場を補正し，最終的な速度場 $\bm{u}^{n+1}$ を得る．
\[ \bm{u}^{n+1}_{i,j} = \bm{u}^*_{i,j} - \frac{\Delta t}{\rho^{n+1}_{i,j}} \text{Grad}_h(p^{n+1})_{i,j} \]
この $\bm{u}^{n+1}$ は，離散的に非圧縮条件を満たす速度場となる．

\end{tcolorbox}
